\clearpage\begingroup\SourceLeaf{415}\ArithmeticHead{84}{L}{STÄCKEL, GAUSS ALS GEOMETER.}
\renewcommand{\thefootnote}{\arabic{footnote})}\setcounter{footnote}{0}
\begin{center}
\textls[100]{Abschnitt V.}\par\smallskip
\textbf{Die allgemeine Lehre von den krummen Flächen.}\par\medskip
\textbf{28.}\par\smallskip
\textls[100]{Entwicklung der Grundgedanken bis zum Jahre 1816.}
\end{center}
\noindent\hspace*{\parindent}Ähnlich wie im 17. Jahrhundert aus den Bedürfnissen der Mechanik die
Infinitesimalrechnung hervorgegangen ist, verdankt im 19. Jahrhundert die
allgemeine Lehre von den krummen Flächen ihre Entstehung der Geodäsie.
In beiden Fällen hat sich aus der angewandten Mathematik ein neuer, lebensfähiger
Zweig der reinen Mathematik losgelöst, hat ein selbständiges Dasein
gewonnen und sich zu einem ausgedehnten, reich gegliederten Inbegriff theoretischer
Untersuchungen ausgestaltet.
\par
Die Frage nach der Gestalt und der Grösse der Erde hatte die Astronomen,
Physiker und Mathematiker während des 18. Jahrhunderts lebhaft beschäftigt,
ja die grossen Gradmessungen in Lappland (1736—1737) und Peru
(1735—1741) hatten die Aufmerksamkeit aller Gebildeten erregt. Handelte
es sich hier um einen rein wissenschaftlichen Gegenstand, so gewann die
Geodäsie bald auch praktische Wichtigkeit. Die Einführung des metrischen
Systems veranlasste die Gradmessung von \textsc{Méchain} und \textsc{Delambre} zwischen
Dünkirchen und Barcelona (1792—1798). Dazu kamen die Anforderungen
der Heeresführung und der Steuerverwaltung, die eine planmässige Triangulierung
der Staaten nötig machten. Hand in Hand mit der Ausdehnung der
geodätischen Messungen ging die Ausbildung und Verfeinerung der mathematischen
Hilfsmittel.
\par
Im Jahre 1816 hatte \textsc{Schumacher}, seit 1815 Leiter der Sternwarte zu
Altona, vom König Friedrich VI. von Dänemark den Auftrag erhalten, Gradmessungen
im Meridian von Skagen bis Lauenburg und im Parallel von
Kopenhagen bis zur Westküste Jütlands als Grundlage für eine spätere Triangulierung
auszuführen. Als \textsc{Schumacher} sogleich bei \textsc{Gauss} anfragte, ob es
sich ermöglichen liesse, den Meridianbogen durch das Königreich Hannover
fortzusetzen und so den Anschluss an die Dreiecke des preussischen
\clearpage\SourceLeaf{416}\ArithmeticHead{85}{R}{DIE ALLGEMEINE LEHRE VON DEN KRUMMEN FLÄCHEN.}\setcounter{footnote}{0}
\noindent Generalstabs zu gewinnen, antwortete dieser am 5. Juli 1816 mit einer bei ihm ungewöhnlichen
Wärme des Tones:
\par
»Vor allen Dingen meinen herzlichen Glückwunsch zu der herrlichen,
grossen Unternehmung, die Sie mir in Ihrem letzten Briefe ankündigen.
Diese Gradmessung in den k. dänischen Staaten wird uns, an sich schon,
über die Gestalt der Erde schöne Aufschlüsse geben. Ich zweifle indessen
gar nicht, dass es in Zukunft möglich zu machen sein wird, Ihre Messungen
durch das Königreich Hannover südlich fortzusetzen.\ .\ .\ .\quad Über die Art, die
gemessenen Dreiecke im Kalkül zu behandeln, habe ich mir eine Methode
entworfen, die aber für einen Brief viel zu weitläufig würde. In Zukunft
.\ .\ .\ werde ich mit Ihnen darüber ausführlich konferieren: ja ich erbiete
mich, die Berechnung der Hauptdreiecke selbst auf mich zu nehmen« (W. IX,
S. 345).
\par
Dass \textsc{Gauss} Freude an geodätischen Messungen und Rechnungen hatte,
lässt sich bis in die Frühzeit hinein verfolgen. Zum Beispiel beteiligte er
sich im August und September 1803 an den Beobachtungen der Pulversignale,
die v. \textsc{Zach} auf dem Brocken veranstaltete, und lieferte um dieselbe Zeit
Berechnungen für die von dem preussischen Generalmajor v. \textsc{Lecoq} vorgenommene
trigonometrische Aufnahme Westfalens\footnote{Näheres hierüber findet man in dem Aufsatz von A. \textsc{Galle} über die geodätischen Arbeiten von \textsc{Gauss}, Werke XI\textsubscript{2}, Abh. 3.}.
\par
Als \textsc{Gauss} im September 1812 v. \textsc{Zach} auf der Sternwarte Seeberg bei
Gotha besuchte, fand er (T. Nr. 142) seine Auflösung der Aufgabe, die Anziehung
eines elliptischen Sphäroids zu bestimmen, die er 1813 veröffentlicht
hat (W. V, S. 1). In der Selbstanzeige sagt er, die Auflösung sei so ausführlich
dargestellt, um sie »auch weniger geübten Lesern verständlich zu
machen, denen diese für die Gestalt der Erde so interessanten Untersuchungen
bisher ganz unzugänglich waren« (W V, S. 217).
\par
Dass \textsc{Gauss} sich in der Zeit zwischen 1812 und 1816 mit der Lehre von
den kürzesten Linien auf dem elliptischen Sphäroid beschäftigt hat, zeigt schon
der vorhin angeführte Brief an \textsc{Schumacher} vom 5. Juli 1816. Dazu kommen
die Briefe an \textsc{Olbers} vom 13. Januar 1821 (W. IX, S. 367) und an \textsc{Bessel}
vom 11. März 1821 und 15. November 1822 (Br. G.-\textsc{Bessel}, S. 380 und 410),
in denen er bemerkt, er habe seine Theorie der Behandlung der Messungen
\clearpage\SourceLeaf{417}\ArithmeticHead{86}{L}{STÄCKEL, GAUSS ALS GEOMETER.}\setcounter{footnote}{0}
\noindent auf der Oberfläche der Erde schon seit geraumer Zeit entwickelt; seine Andeutungen
lassen erkennen, dass er damit das in den artt. 11 und 16 der
\textit{Untersuchungen über Gegenstände der höheren Geodäsie} (W. IV, S. 274 und
286) dargelegte Verfahren meinte. Auch erklärt er in der Selbstanzeige der
zweiten Abhandlung über Gegenstände der höheren Geodäsie vom 28. September
1846: »Der Verfasser, welcher alle diese Untersuchungen schon vor
mehr als dreissig Jahren zu seinem Privatgebrauch durchgeführt und nur bisher
zur Veröffentlichung noch keine Veranlassung gefunden hatte .\ .\ .« (W. IV, S. 353).
\par
\textsc{Gauss} hat jedoch damals noch mehr besessen. Er kannte zunächst die
in dem Brief an \textsc{Schumacher} vom 21. November 1825 (W. VIII, S. 401) erwähnte
Verallgemeinerung des \textsc{Legendre}schen Lehrsatzes von der Zurückführung
eines kleinen sphärischen Dreiecks auf ein ebenes Dreieck mit eben
so langen Seiten. Ferner wird schon in § 10 der \textit{Theoria attractionis corporum
sphaeroidicorum ellipticorum} (1813) auf die Lehre von der Abbildung der krummen
Flächen hingewiesen (W. V, S. 14). Im Frühjahr 1816 hatte \textsc{Gauss} als
Preisaufgabe für die neue, von v. \textsc{Lindenau} und \textsc{Bohnenberger} begründete
»Zeitschrift für Astronomie und verwandte Wissenschaften« die Aufgabe vorgeschlagen,
zwei krumme Flächen mit Erhaltung der Ähnlichkeit in den kleinsten
Teilen auf einander abzubilden\footnote{Hierauf beziehen sich die Briefe von v. \textsc{Lindenau} an \textsc{Gauss} vom 18. und 28. Juni 1816 (Briefe im \textsc{Gauss}-Archiv); die Briefe von \textsc{Gauss} an v. \textsc{Lindenau} scheinen vernichtet worden zu sein, vgl. Br. G.-\textsc{Bolyai}, S. 156 (Brief von \textsc{Sartorius} v. \textsc{Waltershausen} an W. \textsc{Bolyai} vom 12. August 1856).}. Der Brief an \textsc{Schumacher} vom 5. Juli 1816
(W. VIII, S. 370) beweist, dass er ihre Lösung besass; übrigens hat er diese
in einer gleichzeitig niedergeschriebenen Aufzeichnung angegeben (W. VIII,
S. 371). Unmittelbar darauf folgt (Handbuch 16 Bb, S. 71) das »schöne Theorem«,
dass einander entsprechende Stücke von Biegungsflächen, wenn sie auf
die Himmelskugel mittels paralleler Normalen abgebildet werden, auf der
Kugel Flächenstücke gleichen Inhalts ergeben (W. VIII, S. 372). Hierin liegt
die Erhaltung der Gesamtkrümmung eines Flächenstückes gegenüber Biegungen.
Aber auch der Begriff, freilich nicht der Name, des Krümmungsmasses lässt
sich bis in die Zeit zwischen 1813 und 1816 zurückverfolgen, denn eine Notiz
aus dieser Zeit bringt den Satz, dass bei jener Abbildung auf die Kugel vom
Halbmesser Eins das Verhältnis des Bildes eines Flächenelementes zu diesem
selbst gleich dem Produkte der Hauptkrümmungen ist (W. VIII, S. 367).
\clearpage\SourceLeaf{418}\ArithmeticHead{87}{R}{DIE ALLGEMEINE LEHRE VON DEN KRUMMEN FLÄCHEN.}\setcounter{footnote}{0}
\noindent\hspace*{\parindent}Zusammenfassend und in einigen Punkten ergänzend kann man die Ergebnisse
aus der allgemeinen Lehre von den krummen Flächen, zu denen \textsc{Gauss}
bis zum Jahre 1816 gelangt war, etwa folgendermassen darstellen:
\par
1. Auffassung der kartesischen Koordinaten eines Punktes einer krummen
Fläche als Funktionen von zwei Hilfsgrössen (\textit{Theoria attractionis}, § 10, W. V,
S. 14), Abbildung krummer Flächen (ebenda), Abbildung mittels paralleler Normalen
auf die Kugel vom Halbmesser Eins (W. VIII, S. 367)\footnote{Die Beziehung der Richtungen im Raume auf die Punkte der Einheitskugel findet sich schon in der Scheda Ac, Varia, begonnen Nov. 1799, S. 3.}, konforme Abbildung
zweier krummer Flächen auf einander (W. VIII, S. 370).
\par
2. Abwicklung oder Biegung krummer Flächen als besonderer Fall der
Abbildung; Begriff der Gesamtkrümmung eines Flächenstücks, Begriff des
einem Punkte der Fläche zugeordneten Krümmungsmasses, Erhaltung des
Krümmungsmasses gegenüber Biegungen (W. VIII, S. 376, 372).
\par
3. Die Haupteigenschaften der kürzesten Linien auf krummen Flächen,
genauere Untersuchung für das elliptische Sphäroid (W. IX, S. 72—77), Verallgemeinerung
des \textsc{Legendre}schen Theorems auf beliebige Flächen (W. VIII,
S. 401).
\par
Man erkennt, dass bereits in der Zeit zwischen 1812 und 1816 die Fundamente
für das Gebäude der \textit{Disquisitiones generales} gelegt worden sind. Diese
Leistung tritt jedoch erst in das rechte Licht, wenn man sich die gesamte
Tätigkeit von \textsc{Gauss} während jenes Zeitraumes vergegenwärtigt.
\par
In der reinen Mathematik hatte das Jahr 1812 mit der Veröffentlichung
des ersten Teiles der Untersuchungen über die hypergeometrische Reihe begonnen
(W. III, S. 123). Im Dezember 1815 und im Januar 1816 wurden
der Göttinger Gesellschaft die beiden neuen Beweise für den Fundamentalsatz
der Algebra vorgelegt (W. III, S. 31 und 57). Für die Zahlentheorie ist die
im Februar 1817 vorgelegte Abhandlung über die quadratischen Reste zu
nennen, die den fünften und den sechsten Beweis für das Reziprozitätsgesetz
enthält (W. II, S. 47); auch die Lehre von den biquadratischen Resten ist
damals gefördert worden, wie aus den Briefen an \textsc{Bessel} vom 23. Dezember
1816 (W. X\textsubscript{1}, S. 76) und an \textsc{Dirichlet} vom 30. Mai 1828 (W. II, S. 516)
hervorgeht. Aus der Geometrie sind die Untersuchungen zur Flächentheorie
bereits erwähnt worden. Dazu kommen aus dem Jahre 1816 zwei
\clearpage\SourceLeaf{419a}\ArithmeticHead{88}{L}{STÄCKEL, GAUSS ALS GEOMETER.}
\noindent Besprechungen von Versuchen, das Parallelenaxiom zu beweisen (W. IV, S. 363,
VIII, S. 170). Wie wir gesehen haben, wusste \textsc{Gauss} hier mehr, als er öffentlich
auszusprechen für gut fand; er war gerade damals zur nichteuklidischen
Trigonometrie durchgedrungen (W. VIII, S. 176).
\par
Die Abhandlung über die mechanische Quadratur vom 16. September
1814 bildet den Übergang zur angewandten Mathematik (W. III, S. 163). In
diese selbst gehört die Bestimmung der Anziehung der homogenen elliptischen
Sphäroide vom 18. März 1813 (W. V, S. 1). Im Anschluss an die Beobachtungen
des Kometen vom Jahre 1813 wurde die \textit{Theoria motus} nach der
Seite der parabolischen Bahnen ergänzt; die betreffende Abhandlung ist vorgelegt
am 10. September 1813 (W. VI, S. 25). Ferner sind anzuführen zahlreiche,
meistens in den Göttinger Anzeigen veröffentlichte astronomische Rechnungen
und Beobachtungen (W. VI, S. 354—392). Die Untersuchungen aber,
denen \textsc{Gauss} während der Zeit von 1810 bis 1818 wohl den grössten Teil
seiner Zeit und Kraft gewidmet hat, die Störungen der Pallas, sind nicht abgeschlossen
worden; erst im Jahre 1906 hat \textsc{Brendel} die Bruchstücke herausgegeben
(W. VII, S. 439—600).
\par
»In jener Zeit« schreibt \textsc{Sartorius von Waltershausen} (S. 50), »schien
ihm keine Anstrengung des Geistes und des Körpers zu gross, um eine Reihe
von Arbeiten durchzuführen, dazu bestimmt, die Wissenschaft des 19. Jahrhunderts
zu reformieren und ihr Fundamente zu unterbreiten, deren Festigkeit
erst von künftigen Geschlechtern anerkannt und gewürdigt werden wird«.
\clearpage\endgroup
